\chapter{JavaScript: prototipi, classi ed ereditariet\`a}

\section{I prototipi}

\begin{definizione}[Prototipo]
In JavaScript \textbf{ogni oggetto ha un prototipo}: un altro oggetto di cui \emph{eredita le
propriet\`a}. L'unico oggetto \emph{privo} di prototipo \`e \texttt{Object.prototype}, che \`e
il prototipo di default degli object literal.
\end{definizione}

Due oggetti ``primordiali'' formano la base di tutte le strutture dati del linguaggio:
\begin{itemize}
  \item la \textbf{funzione \texttt{Object}}: invocata implicitamente ogni volta che c'\`e da
        creare un oggetto (es. dichiarando un object literal);
  \item l'oggetto \textbf{\texttt{Object.prototype}}: prototipo di default per tutti gli
        oggetti creati senza specificarne uno diverso. (La scrittura non deve stupire: in JS
        una funzione \emph{\`e} un oggetto, e pu\`o quindi possedere propriet\`a.)
\end{itemize}

\begin{importante}[Come funziona l'ereditariet\`a via prototipo]
Quando si \emph{legge} una propriet\`a di un oggetto (per conoscerne il valore o per invocare
una propriet\`a-funzione), se il motore non la trova nell'oggetto stesso la cerca \textbf{nel
suo prototipo}: per questo si dice che l'oggetto \emph{eredita} le propriet\`a del prototipo.
Il prototipo viene invece \textbf{ignorato in scrittura}: assegnando un valore a una
propriet\`a inesistente, questa viene \emph{aggiunta all'oggetto}. Se il prototipo aveva una
propriet\`a con lo stesso nome siamo in presenza di \textbf{overriding}: l'erede sovrascrive
(per s\'e) la propriet\`a dell'antenato -- il motore trover\`a prima quella ``locale'' e
ignorer\`a quella ereditata. Gli \emph{altri} eredi continuano a condividere quella del
prototipo.
\end{importante}

\begin{definizione}[Prototype chain]
Ogni oggetto ha un prototipo, che \`e a sua volta un oggetto con un prototipo, e cos\`i via
fino a \texttt{Object.prototype}: questa \`e la \textbf{catena dei prototipi}
(prototype chain). Prima di decidere che un oggetto \emph{non} possiede una propriet\`a, il
motore risale \emph{tutta} la catena.
\end{definizione}

\subsection{Accedere e modificare i prototipi}
Si pu\`o accedere in ogni momento al prototipo di un oggetto tramite
\texttt{oggetto.\_\_proto\_\_}, e attraverso di esso \emph{modificarne le propriet\`a} --
di fatto modificando ci\`o che \`e accessibile a \textbf{tutti} gli oggetti che fanno capo a
quel prototipo (eliminando \texttt{abbaia} dal prototipo \texttt{fido}, anche
\texttt{snoopy} e \texttt{cucciolo} la perdono). Vale anche per i prototipi predefiniti:
aggiungendo una propriet\`a a \texttt{Object.prototype}, \emph{tutti} gli oggetti
dell'applicazione la acquisiscono -- si pu\`o ``plasmare o distruggere il linguaggio'' a
proprio rischio e pericolo.

\`E possibile (ma poco usato) bloccare le modifiche a un oggetto dopo la creazione:
\begin{itemize}
  \item \texttt{Object.preventExtensions(ogg)}: non potr\`a acquisire nuove propriet\`a n\'e
        cambiare prototipo;
  \item \texttt{Object.seal(ogg)}: \emph{sigilla} -- come sopra, ma non potr\`a nemmeno
        \emph{perdere} propriet\`a;
  \item \texttt{Object.freeze(ogg)}: \emph{congela} -- oltre a sigillare, impedisce nuovi
        assegnamenti alle propriet\`a: di fatto read-only.
\end{itemize}

\subsection{Creare oggetti a partire da un prototipo}
\begin{lstlisting}[language=JavaScript]
const fido = {              // prototipo di default: Object.prototype
  tipo: "Cane",
  nome: "Fido",
  abbaia: function () {}
}
const snoopy = Object.create(fido);   // fido come prototipo
snoopy.nome = "Snoopy";               // overriding di "nome"
snoopy.razza = "Beagle";              // proprieta' originali
const cucciolo = Object.create(fido);
cucciolo.razza = "Jack Russel Terrier";
// cucciolo.nome vale "Fido" finche' non lo ridefinisce
\end{lstlisting}

\subsection{Determinare le propriet\`a di un oggetto}
\begin{itemize}
  \item \texttt{Object.keys(ogg)}: array dei nomi delle propriet\`a definite
        \emph{nell'oggetto stesso};
  \item \texttt{for (let key in ogg)}: cicla su \emph{tutte} le chiavi valide, \textbf{incluse
        quelle della prototype chain} ma \emph{escluse} quelle ``di sistema'' (es. le
        propriet\`a definite in \texttt{Object.prototype});
  \item \texttt{ogg.hasOwnProperty("p")}: metodo definito in \texttt{Object.prototype} (dunque
        disponibile a tutti): true se l'oggetto \emph{definisce} la propriet\`a; false se non
        esiste \emph{o \`e ereditaria}.
\end{itemize}

\section{Metodi e binding di \texttt{this}}

\begin{definizione}[Metodo]
Un \textbf{metodo} \`e una funzione \emph{collegata} (bound) a un oggetto, che diventa il
``soggetto'' della funzione: nell'esecuzione, \texttt{this} si riferisce a
quell'oggetto-soggetto.
\end{definizione}

\begin{importante}[Quando avviene il binding]
\begin{itemize}
  \item Il binding avviene \textbf{implicitamente} per le propriet\`a-funzione definite come
        \emph{named} o \emph{anonymous} function.
  \item \textbf{Non avviene} per le propriet\`a-funzione definite come \emph{arrow function}:
        una arrow assegnata a una propriet\`a \emph{non} \`e un metodo (dentro di essa
        \texttt{this} non \`e l'oggetto proprietario).
  \item \textbf{Binding esplicito}: \texttt{f.bind(ogg)} crea una \emph{copia} della funzione
        \texttt{f} bound all'oggetto specificato. Lo stesso effetto si ottiene assegnando la
        funzione a una propriet\`a dell'oggetto -- a meno che non sia una arrow, nel qual caso
        il binding potr\`a essere solo esplicito.
\end{itemize}
\begin{lstlisting}[language=JavaScript]
const fido = {
  tipo: "Cane", nome: "Fido",
  abbaia: function () {},  // metodo (binding implicito)
  corre: () => {}          // NON metodo (arrow: niente binding)
}
function log(msg) { console.log(msg + ": " + this.nome); }
const stampaFido = log.bind(fido);  // binding esplicito
fido.stampa = log;                   // binding implicito
\end{lstlisting}
\end{importante}

\section{I costruttori: creare oggetti ``con lo stampino''}

Il sistema originario per creare oggetti simili in serie sono i \textbf{costruttori}:
\emph{named function} (per convenzione con l'iniziale maiuscola) invocate premettendo la
keyword \textbf{new}.
\begin{lstlisting}[language=JavaScript]
function Cane(nome, razza) {
  this.tipo = "Cane";
  this.nome = nome;
  if (razza !== undefined) this.razza = razza;
}
const fido = new Cane("Fido");
const snoopy = new Cane("Snoopy", "Beagle");
\end{lstlisting}

\begin{importante}[Cosa succede con \texttt{new Fun(\dots)}]
\begin{enumerate}
  \item viene creato un oggetto da usare come \emph{prototipo}, memorizzato in
        \texttt{Fun.prototype};
  \item viene creato un nuovo oggetto \texttt{x} con prototipo \texttt{Fun.prototype};
  \item viene realizzato il \emph{binding} fra il nuovo oggetto e la funzione: \texttt{this}
        dentro la funzione si riferisce a \texttt{x}.
\end{enumerate}
Dunque \textbf{tutti gli oggetti creati con \texttt{new Fun(\dots)} condividono il medesimo
prototipo}.
\end{importante}

\section{Classi = costruttore + prototipo condiviso}

\begin{definizione}[Classe in JavaScript]
Il concetto di costruttore \`e alla base del concetto di classe: \emph{due oggetti sono della
stessa classe se creati dallo stesso costruttore}. Oggetti creati dallo stesso costruttore
sono inizializzati allo stesso modo (stesse propriet\`a, ognuno coi propri valori: la
propriet\`a \`e \emph{dell'oggetto}, non del prototipo) e \textbf{condividono il prototipo}:
per questo si possono dare ``metodi alla classe'' semplicemente aggiungendoli al prototipo.
\end{definizione}

\begin{esempio}[Il costrutto class \`e zucchero sintattico]
Le due forme seguenti sono equivalenti: la ``classe'' non \`e altro che un costruttore, e i
metodi scritti ``dentro'' la classe vengono assegnati al prototipo associato al costruttore.
\begin{lstlisting}[language=JavaScript]
// con costruttore e prototipo esplicito
function Cane(nome, razza) {
  this.tipo = "Cane"; this.nome = nome;
  if (razza !== undefined) this.razza = razza;
}
Cane.prototype.abbaia = function() {}

// con il costrutto class
class Cane {
  constructor(nome, razza) {
    this.tipo = "Cane"; this.nome = nome;
    if (razza !== undefined) this.razza = razza;
  }
  abbaia() {}
}
\end{lstlisting}
\end{esempio}

\subsection{Ereditariet\`a fra classi e catena dei prototipi}
Con \texttt{class Beagle extends Cane}, la catena che si forma \`e:
\begin{center}
\texttt{snoopy} $\rightarrow$ \texttt{Beagle.prototype} $\rightarrow$
\texttt{Cane.prototype} $\rightarrow$ \texttt{Object.prototype}
\end{center}
\begin{lstlisting}[language=JavaScript]
class Cane {
  constructor(nome, razza) { ... }
  abbaia() {}
}
class Beagle extends Cane {
  constructor(nome) {
    super("Beagle");   // prima cosa: costruttore della superclasse
    ...
  }
  scrive() {}
}
const snoopy = new Beagle("Snoopy");
// snoopy trova scrive() in Beagle.prototype,
// abbaia() risalendo a Cane.prototype
\end{lstlisting}
